% !PRINTED 572
\weberpage{572}
Dass diese Formen aber wirklich von einander verschieden sind,
ergiebt sich daraus, dass wir (nach \S.\,145, 150) für die Variablen
$t_1, t_2, \ldots, t_n$ solche ganze rationale Zahlen setzen können,
dass $\tau$ in eine primitive Wurzel $\gamma$ des Primideals
$\mathfrak p$ übergeht. Durch dieselbe Substitution werden die
Grössen (10)
\weq{11}{\equiv \gamma, \gamma^p, \gamma^{p^2}, \ldots, \gamma^{p^{f-1}}\ (\mathrm{mod}\ \mathfrak p),}
die nach dem Modul $\mathfrak p$ alle von einander verschieden sind.
Es können also auch nicht zwei der Formen (10) nach dem Modul
$\mathfrak p$ congruent sein, weil sonst auch die beiden
entsprechenden Zahlen (11) congruent ausfallen würden.

Dies fassen wir als Satz so zusammen:

1. \gesperrt{Jede Congruenz (4), deren eine Wurzel $t = \tau$ ist,
hat die $f$ verschiedenen Wurzeln}
\[
\tau, \tau_1, \tau_2, \ldots, \tau_{f-1}.
\]

Hiernach können wir, indem wir $\Phi(t)$ durch $t - \tau$ algebraisch
dividiren,
\[
\Phi(t) \equiv (t - \tau)\, \Phi_1(t)\ (\mathrm{mod}\ \mathfrak p)
\]
setzen, und $\tau_1, \tau_2, \ldots, \tau_{f-1}$ sind Wurzeln von
$\Phi_1(t) \equiv 0$, worin aber $\Phi_1(t)$ noch nicht rationale
Coëfficienten, sondern Coëfficienten in $\mathfrak o$ hat. Dividiren
wir $\Phi_1(t)$ wieder durch $t - \tau_1$, und fahren so fort, so
folgt endlich, wenn wir also nun die Function $f^{\text{ten}}$ Grades
\weq{12}{\Pi(t) = (t - \tau)(t - \tau_1) \ldots (t - \tau_{f-1})}
setzen,
\weq{13}{\Phi(t) \equiv \Pi(t)\, \Phi_0(t)\ (\mathrm{mod}\ \mathfrak p),}
worin $\Phi_0(t)$ eine ganze Function mit Coëfficienten in
$\mathfrak o$ ist.

Die Function $\Pi(t)$ hängt von den Variablen $t, t_1, \ldots, t_n$
ab, und um dies auszudrücken, setzen wir
\[
\Pi(t) = \Pi(t, t_1, \ldots, t_n).
\]

Die Coëfficienten dieser Form sind Zahlen in $\mathfrak o$, und es
lässt sich noch nachweisen, dass sie mit ganzen rationalen Zahlen
nach dem Modul $\mathfrak p$ congruent sind. Dieser Beweis ergiebt
sich durch Erheben in die $p^{\text{te}}$ Potenz:
\[
[\Pi(t)]^p \equiv (t^p - \tau^p)(t^p - \tau_1^p) \ldots (t^p - \tau_{f-1}^p).
\]

Nun ist aber nach (9)
\[
\tau_r^p \equiv \omega_1^{p^{r+1}} t_1^p + \omega_2^{p^{r+1}} t_2^p + \cdots + \omega_n^{p^{r+1}} t_n^p\ (\mathrm{mod}\ \mathfrak p),
\]
und wir erhalten also $\tau_r^p$ aus $\tau_{r+1}$, wenn wir $t_i$
durch $t_i^p$ ersetzen,
%CONT>>
